\chapter{生产者、broker 与消费者的一致性}

本章沿一条消息查看生产确认、broker 复制和 consumer 处理三个彼此独立的边界。\textbf{结果未知}表示服务端可能已完成动作但回复没到；重试是再次提交同一意图，可能产生重复。\textbf{幂等}表示相同稳定身份与内容的重复请求不会再增加一次持久效果；它不等于事务。至少一次处理愿意重放不确定工作，因此副作用可能重复。\textbf{副作用}是回调对日志以外状态的更改，例如更新本地账本。\textbf{事务}应把多个参与者的更改作为一个不可分割的提交；本课程 producer、复制日志、consumer offset 和账本之间没有共同事务。exactly-once 若指端到端每个业务事件恰有一次效果，需要协调这些边界；单独的 ACK、HW、commit 或 producer 幂等都不提供该保证。轨迹使用真实 TCP 和各 broker 独立日志，分开记录物理 LEO、authority 报告、HW、consumer position、持久 committed offset、可见记录与回调效果。

\lessonsection{29}{三层一致性故障实验}
\goals{通过可复现的 producer、broker 与 consumer 故障，记录各层真实发生的状态。}{从基础区分成功追加、回复丢失、HW 以下可见性、内存 position、group committed offset 与外部副作用。}
\background{先修\stepref{27}{元数据路由刷新}与\stepref{28}{副本重新入ISR}。LEO 是下一条可写的分区 offset；physical LEO 是该 broker 磁盘末端，reported LEO 是 authority 收到的复制进度。HW 是复制提交边界，consumer 只能看到 offset<HW。poll 可能使内存 position 前进，commit 才保存下一次 resume 的位点；两者都不删除日志。callback 可在 commit 前改变外部状态。receipt 是 producer 对追加请求的返回，不是 consumer 可见证明。}
\providededit{provided 层的真实状态工具位于 \pathcode{provided/src/main/java/io/simplekafka/lab/}：\pathcode{FaultProxy.java}、\pathcode{EffectRecorder.java}、\pathcode{ConsistencyObservation.java}、\pathcode{ConsistencyTrace.java}、\pathcode{LabSupport.java}；\pathcode{ClusterHarness} 位于 \pathcode{provided/src/main/java/io/simplekafka/cluster/ClusterHarness.java}。}{\pathcode{exercises/src/main/java/io/simplekafka/lab/ConsistencyExperiments.java}：五个 public static 方法都返回 \pathcode{ConsistencyTrace}：\method{producerOutcomeUnknown(ClusterHarness,TopicPartition,RecordData,FaultProxy)}、\method{leaderAckLoss(ClusterHarness,TopicPartition,RecordData)}、\method{allAckTimeout(ClusterHarness,TopicPartition,RecordData)}、\method{consumerReplay(ClusterHarness,TopicPartition,String,RecordData,EffectRecorder,FaultProxy)}（\pathcode{throws Exception}）与 \method{staleGroupSideEffect(ClusterHarness,TopicPartition,String,RecordData,EffectRecorder)}。}
\paragraph{先理解观察工具。}\pathcode{ConsistencyObservation} 的完整类型为 \pathcode{record ConsistencyObservation(String phase,ErrorCode error,ReplicationSnapshot replication,Map<Integer,Long> physicalLeo,OptionalLong position,OptionalLong committedOffset,List<LogRecord> visible,List<LogRecord> effects,Optional<AppendResult> receipt)}；\pathcode{ConsistencyTrace} 为 \pathcode{record ConsistencyTrace(TopicPartition tp,List<ConsistencyObservation> observations)}。effects 保留每次调用，不去重。提供的 \pathcode{FaultProxy(Endpoint)} 暴露 \pathcode{upstream():Endpoint}、\pathcode{start():Endpoint}、\pathcode{endpoint():Endpoint}、\pathcode{armed():boolean}、\pathcode{dropNextReply(short)}、\pathcode{rejectNextRequest(short)}、\pathcode{forwarded/dropped/rejected(short):int} 和 \pathcode{close()}；同一时刻只准 arm 一个 API。drop 会把请求转给真实 broker 并读到回复后丢弃回复；reject 在转发前返回 \method{REQUEST_TIMEOUT} 与 \pathcode{ErrorBody("injected pre-forward failure")}。不匹配的 API 原样转发。
\pathcode{LabSupport.replicateAndReport(ClusterHarness,TopicPartition):void} 实际通过 TCP 复制并只报告持久后的 LEO；\pathcode{observe(String,ErrorCode,ClusterHarness,TopicPartition,OptionalLong,String,EffectRecorder,Optional<AppendResult>):ConsistencyObservation} 独立读取可见记录与 committed offset；\pathcode{awaitLeaderLeo(ClusterHarness,TopicPartition,long):void} 等待 authority 收到指定进度。provided \pathcode{EffectRecorder.process(TopicPartition,LogRecord):void} 和 \pathcode{snapshot():List<LogRecord>} 按调用顺序记录回调。
\lessoncontract{五个方法都要求调用方提供已启动的 cluster、唯一 partition \pathcode{tp}、非空且未 stamped 的输入记录。初态必须是 epoch0、leader1、replicas/ISR={1,2,3}、三者在线且日志 start=end=0、HW=0、minISR=2。consumer 场景还要求 group 无 committed offset 且 recorder 初始为空；group 与 recorder 必须同时提供或同时不提供。两个 proxy 场景要求调用方持有已启动、尚未 armed 且 upstream 是 broker1 的 proxy。初态不符、参数空值或输入带 stamp 为 \method{INVALID_REQUEST}，且不可追加/改 offset/改 recorder/proxy 故障计数。只关闭方法自有 client；只捕获本实验预期错误，其余异常传播。}
\paragraph{实验一：producer 回复丢失后显式重发。}测试记录取 \pathcode{key="customer-A"}（10 ASCII bytes）、\pathcode{value="request-unknown-73"}（18 bytes）、timestamp=1723；普通记录占 \(32+10+18=60\) bytes。reader 从 position0 开始，producer 设置 \pathcode{batchRecords=1}。时间线中 physical LEO 顺序为 broker 1/2/3；replication 值为 leader/epoch/ISR/reported LEO/HW。
\begin{itemize}
\item \pathcode{UNKNOWN}: broker 已追加 offset0；LEADER 回复丢失，因此 producer 报 \method{REQUEST_TIMEOUT}。\newline 状态：replication=(1/0/{1,2,3}/{1,0,0}/0)，physical LEO=(1,0,0)；\newline position=0，committed=none，visible=[]，effects=[]，receipt=none。客户端没有收到回复。
\item \pathcode{POISONED}: 同一普通 producer 本地拒绝后续 send，error=\method{REQUEST_TIMEOUT}；replication/LEO 仍为上行状态，position=0，committed=none，visible=[]，effects=[]，receipt=none；此时 forwarded(PRODUCE)=1、dropped(PRODUCE)=1，第二次调用没有经过网络。
\item \pathcode{EXPLICIT_RESEND}: 新 producer 显式发送相同数据，得到 receipt=[1,2)。三份日志完成复制并报告后：\newline replication=(1/0/{1,2,3}/{2,2,2}/2)，physical LEO=(2,2,2)；\newline reader position=2，committed=none，visible=[offset0,offset1]（内容相同），effects=[]，error=\method{NONE}。\newline 最终 forwarded(PRODUCE)=2、dropped(PRODUCE)=1。
\end{itemize}
\paragraph{实验二：LEADER receipt 不保护未提交尾部。}测试输入 \pathcode{key="leader-tail"}（11 bytes）、\pathcode{value="acked-but-uncommitted"}（21 bytes）、timestamp=2419，记录64 bytes。启动无 committed offset 的 reader，position=0。
\begin{itemize}
\item \pathcode{ACKED_NOT_VISIBLE}: broker1 追加 offset0 后返回 LEADER receipt=[0,1)，error=\method{NONE}；replication=(1/0/{1,2,3}/{1,0,0}/0)，physical=(1,0,0)，position=0、committed=none、visible=[]、effects=[]。
\item \pathcode{ELECTED_WITHOUT_TAIL}: 停 broker1 后选举 broker2/epoch1；replication=(2/1/{2}/{0,0,0}/0)，physical=(1,0,0)，reader position=0、committed=none、visible=[]、effects=[]、receipt=none、error=\method{NONE}。旧日志还在 broker1 磁盘，但不是新 leader 可见的 HW 前缀。
\item \pathcode{OLD_TAIL_REMOVED}: 重启 broker1，先 \method{reconcile} 再 \method{admit}；offset0 被截去。\newline 新状态：replication=(2/1/{1,2}/{0,0,0}/0)，physical=(0,0,0)；\newline position=0，committed=none，visible=[]，effects=[]，receipt=none，error=\method{NONE}。第一次 LEADER receipt 不会被追溯撤销。
\end{itemize}
\paragraph{实验三：ALL 超时后稍晚可见。}测试记录为 \pathcode{key="all-key"}（7 bytes）、\pathcode{value="late-visible-value"}（18 bytes）、timestamp=3101，占57 bytes；consumer position 初始0。
\begin{itemize}
\item \pathcode{TIMEOUT_NOT_VISIBLE}: 以 \pathcode{Acks.ALL,timeoutMillis=0} 追加，error=\method{REQUEST_TIMEOUT}；replication=(1/0/{1,2,3}/{1,0,0}/0)，physical=(1,0,0)，position=0、committed=none、visible=[]、effects=[]、receipt=none。追加不回滚。
\item \pathcode{LATE_VISIBLE}: 后续完成真实复制并报告三份副本。\newline 新状态：replication=(1/0/{1,2,3}/{1,1,1}/1)，physical=(1,1,1)；\newline poll 后 position=1，committed=none，visible=[offset0]，effects=[]，receipt=none，error=\method{NONE}。迟到的 HW 不会补发或创造原请求 receipt。
\end{itemize}
\paragraph{实验四：commit 失败后副作用重放。}测试记录 \pathcode{key="order-87"}（8 bytes）、\pathcode{value="effect-after-poll"}（17 bytes）、timestamp=4003，占57 bytes。起初 Kafka 日志已经复制到 offset0，HW=1；consumer group 尚无 commit。proxy 在第一次 \pathcode{COMMIT_OFFSET} 转发前拒绝。
\begin{itemize}
\item \pathcode{COMMIT_FAILED}: first consumer resume=0，poll/callback 后 position=1，error=\method{REQUEST_TIMEOUT}。\newline 状态：replication=(1/0/{1,2,3}/{1,1,1}/1)，physical=(1,1,1)；\newline committed=none，visible=[offset0]，effects=[0]，receipt=none。proxy counts: rejected(\pathcode{COMMIT_OFFSET})=1、forwarded(\pathcode{COMMIT_OFFSET})=0.
\item \pathcode{REPLAY_COMMITTED}: 新 consumer 因没有持久位点再次 resume=0，重放 callback 后 commit nextOffset=1 成功；error=\method{NONE}，replication=(1/0/{1,2,3}/{1,1,1}/1)，physical=(1,1,1)，position=1，committed=1，visible=[offset0]，effects=[0,0]，receipt=none。重复的是外部处理，不是 Kafka 日志追加。
\end{itemize}
\paragraph{实验五：stale group side effect。}测试记录 \pathcode{key="inventory-13"}（12 bytes）、\pathcode{value="old-owner-effect"}（16 bytes）、timestamp=5007，占60 bytes；消息已复制并在offset0可见。
\begin{itemize}
\item 成员z以 generation1 取得 partition、fetch offset0 并完成 callback。成员a加入后 generation2 把 partition 分给a。z 用旧 token 的 commit 与 fetch 均得到 \method{ILLEGAL_GENERATION}；观察 phase=\pathcode{STALE_COMMIT_REJECTED}、error=\method{ILLEGAL_GENERATION}、replication=(1/0/{1,2,3}/{1,1,1}/1)、physical=(1,1,1)、position=1、committed=none、visible=[offset0]、effects=[0]、receipt=none。
\item a 从无已提交位点处 fetch 同一 offset0，再处理并提交 nextOffset1。phase=\pathcode{NEW_OWNER_COMMITTED}、error=\method{NONE}，replication/LEO/HW不变，position=1、committed=1、visible=[offset0]、effects=[0,0]、receipt=none。generation fencing 拒绝旧 broker 请求，不能撤销z已经做的 callback。
\end{itemize}
\lessonhints{逐个区别 physical LEO、authority reported LEO、HW、client position 与 durable committed nextOffset；它们是不同状态。}{给五条故障时间线标注 append、复制、callback、commit 与错误发生的先后，不把一个层次的成功推断为另一层次成功。}{对照每阶段真实 snapshot/trace，再独立检查调用方仍持有的 FaultProxy 计数和 EffectRecorder 列表；不以 sleep 制造竞态。}
\expected{先用 \pathcode{UNKNOWN→POISONED→EXPLICIT_RESEND} 检查一次实际 append、一次本地拒绝、一次显式重发；再分别重演 leader election、ALL 超时、commit 拒绝和 generation 改变。注意：LEADER receipt=[0,1)可与 HW0/不可见并存；commit拒绝可与callback已执行并存；group当前 owner 可再次执行该副作用。通过的是状态演算，不是“所有阶段都成功”的假设。}
\paragraph{测试与完成标准。}\method{Step29Test.producerOutcomeUnknownAndExplicitResendAreObserved} 核对 reply drop、poison 与重复 offset；\method{leaderAcknowledgmentDoesNotProtectAnUncommittedTail} 核对 clean election、旧尾删除及 admission；\method{allAcknowledgmentTimeoutCanBecomeVisibleLater} 核对无 receipt 的 late visibility；\method{consumerCommitFailureReplaysEffectAndCommitsOnFreshConnection} 核对 commit 转发计数、position、offset 与 callback 重放；\method{staleMemberSideEffectSurvivesGenerationFencing} 核对旧 token 两个错误与旧副作用。另有 proxy 单故障、stamped input、非空日志和无效初态不改变状态的测试。单步命令定位Step29；累计命令运行Steps1–29。
\pitfalls{receipt 只说明某种 ACK 的结果；LEADER 已确认的尾部仍可在选举中丢失。timeout不证明未追加。position 在 commit 前推进；generation fencing 约束 broker 请求而非外部 callback。这些实验不重试普通 \method{SimpleProducer}，不声称端到端事务或 exactly-once。}
\lessoncommand{29}
\lessonquestions{哪些观察能区分 LEADER 确认与复制提交后的可见性？}{为什么旧成员 commit 被拒，仍可能已有副作用？}
\paragraph{自检答案。}检查 receipt、physical/reported LEO、HW 与 visible：LEADER receipt 可先于复制和可见。callback 与 commit 是两个动作，commit 失败只说明进度未持久化；它不会回滚已经发生的效果。

\lessonsection{30}{幂等生产者}
\goals{持久保存 producer batch 身份，使同一批次的显式重试返回原 offset 区间而不重复追加。}{从基本身份概念出发，贯通 stamped record、日志恢复、复制 ACK、路由和显式 client retry。}
\background{先修\stepref{29}{三层故障实验}以及\stepref{1}{记录格式}、\stepref{6}{分段日志}与\stepref{24}{复制写入}。网络错误可能让 producer 不知道 broker 是否追加；若没有持久身份，直接重发就是一条新记录。幂等追加给一批逻辑记录稳定身份，broker 可以把完全相同的重试映射回第一次的 offset。四种编号不可混淆：leader epoch 是 broker 角色任期；producer epoch 是调用方的 producer 会话任期；firstSequence 是该 producer 在该分区中下一个 batch 序号；Kafka offset 是分区日志中的位置。eventId 则标识业务事件，第31步用它跨不同 offset 去重。}
\providededit{provided 层提供 \pathcode{ProducerStamp}（\pathcode{provided/src/main/java/io/simplekafka/model/ProducerStamp.java}）与 \pathcode{RecordData}（\pathcode{provided/src/main/java/io/simplekafka/model/RecordData.java}），\pathcode{ProducerBatchFingerprint.fingerprint(List<RecordData>):byte[]}（\pathcode{provided/src/main/java/io/simplekafka/protocol/ProducerBatchFingerprint.java}）、\pathcode{ProducerBatchSupport.validateAppend(ProducerStamp,List<RecordData>):void}（\pathcode{provided/src/main/java/io/simplekafka/storage/ProducerBatchSupport.java}）、\pathcode{Messages.IdempotentProduceRequest}、\pathcode{MessageCodec.encodeRequest(Messages.Request):byte[]}、\pathcode{decodeRequest(short,byte[]):Messages.Request}（\pathcode{provided/src/main/java/io/simplekafka/protocol/Messages.java}、\pathcode{provided/src/main/java/io/simplekafka/protocol/MessageCodec.java}）、以及 \pathcode{IdempotentProduceBackend.produceIdempotent(List<RecordData>,Acks,int,long,long,int,long):AppendResult}（\pathcode{provided/src/main/java/io/simplekafka/broker/IdempotentProduceBackend.java}）。\method{IdempotentProducer(MetadataRouter,long,int)} 构造器/\method{close()} 与 \method{IdempotentAppender(PartitionLog)} 构造器在对应 exercise 类中已给，Appender 不会自动 recover。}{按依赖次序完成六个子任务；实际编辑路径与方法签名如下。}
\begin{enumerate}
\item \textbf{身份与格式。} provided 的 \pathcode{ProducerStamp(long producerId,int producerEpoch,long firstSequence,int batchSize,int batchIndex,byte[] batchHash)} 对每条记录保存生产者 ID、producer epoch、batch 首序号、批量大小、零起批内索引及32-byte SHA-256 hash。id/epoch/sequence非负、batchSize>0、index属于[0,batchSize)、sequence+size不溢出。hash输入为大端 \pathcode{count:int32}，逐记录 \pathcode{timestamp:int64 | keyLength:int32(-1=null,0=empty) | key | valueLength:int32 | value}；不包括offset/leader epoch。普通记录字节维持原格式；stamped记录多64 bytes。磁盘格式为 \pathcode{length | crc32c | offset | timestamp | marker:int32=-2 | producerId:int64 | producerEpoch:int32 | firstSequence:int64 | batchSize:int32 | batchIndex:int32 | batchHash:32 | keyLength:int32 | valueLength:int32 | key | value}，CRC 覆盖 offset 至 value。普通记录最小总长32、stamped最小96；length字段上限1,048,576、总编码最多1,048,580 bytes（4-byte前缀另计入总长）。null key 时 stamped value 最大1,048,484 bytes，key/value均计入限制。教学 wire 的stamp布局不同于磁盘：marker、identity、hash、keyLength/key、valueLength/value、timestamp。
\item \textbf{记录编解码与单段日志。} 在 \pathcode{exercises/src/main/java/io/simplekafka/storage/RecordCodec.java} 扩展 \method{public static ByteBuffer encode(LogRecord)} 和 \method{public static LogRecord decode(ByteBuffer)}，保留普通记录格式/CRC/原有 buffer 契约，并以 marker区分 stamped 记录。扩展 \pathcode{exercises/src/main/java/io/simplekafka/storage/SegmentLog.java} 的 \method{public AppendResult append(List<RecordData>)} 与 \method{public long recover()}：写入、恢复时维护/重建尾 ProducerStamp，验证记录结构与CRC。provided helper \method{void truncateToOffset(long)} 只是 SegmentLog 的截尾操作，不是 PartitionLog 的新 API。
\item \textbf{分区日志与前缀版本。} 在 \pathcode{exercises/src/main/java/io/simplekafka/storage/PartitionLog.java} 扩展 \method{public synchronized AppendResult append(List<RecordData>)}、\method{public synchronized long deleteBefore(long offset)}、\method{public synchronized void truncateTo(long nextOffset)}；提供 \method{long mutationVersion()}、\method{ProducerStamp tailProducerStamp()}、\method{long prefixVersion()}。append 要求非空、合法 batch，连续分配offset，整批校验后按段追加并更新尾stamp；tail没有 stamped 记录时 \pathcode{tailProducerStamp()} 为null。deleteBefore只删 \pathcode{endOffset<=offset} 的非active封闭段、至少保留一个段，负offset为 \method{INVALID_REQUEST}，返回实际 logStart；truncateTo 的合法域是 [logStart,LEO]，越界为 \method{OFFSET_OUT_OF_RANGE}，等于LEO不变。有效日志mutation在首次I/O前记一次；实际前缀改写/保留删除增加prefixVersion，no-op不使缓存失效。I/O失败可能已改磁盘，错误状态按日志的 failed/reopen 约定处理。不可把 SegmentLog 的 \method{truncateToOffset(long)} 写成 PartitionLog API。
\item \textbf{日志派生的 batch 去重。} 在 \pathcode{exercises/src/main/java/io/simplekafka/replication/IdempotentAppender.java} 实现 \method{public AppendResult append(long producerId,int producerEpoch,long firstSequence,List<RecordData> records)} 和 \method{public void recover()}。构造器 \method{IdempotentAppender(PartitionLog)} 已给。append 接收普通 unstamped 输入，验证完整 batch 后才追加；recover 从保留日志派生 producer epoch、nextSequence、完成批次范围与hash，不把易失 cache 当持久真相。日志start>0时报 \method{PRODUCER_STATE_EXPIRED}；完整历史格式/序列/hash损坏报 \method{CORRUPT_RECORD}。
\item \textbf{复制写入与 API 11。} 扩展 \pathcode{exercises/src/main/java/io/simplekafka/replication/ReplicatedPartition.java} 的 \method{public AppendResult produceIdempotent(List<RecordData> records,Acks acks,int leaderEpoch,long timeoutMillis,long producerId,int producerEpoch,long firstSequence)}；扩展 \pathcode{exercises/src/main/java/io/simplekafka/broker/BrokerHandler.java} 的 \method{public Messages.Reply handle(Messages.Request request)} 分派 API \pathcode{IDEMPOTENT_PRODUCE=11}，成功仍回 \pathcode{ProduceBody(AppendResult)}；扩展 \pathcode{exercises/src/main/java/io/simplekafka/client/MetadataRouter.java} 的 \method{public Messages.Reply call(TopicPartition tp,Messages.Request request)} 路由 API11。provided 的请求记录签名为 \pathcode{IdempotentProduceRequest(TopicPartition tp,int epoch,Acks acks,long timeoutMillis,long producerId,int producerEpoch,long firstSequence,List<RecordData> records)}，record输入按序且未 stamped。单 broker backend 不支持该 capability，须返回 \method{INVALID_REQUEST} 且不得追加。
\item \textbf{client 与显式 retry。} 在 \pathcode{exercises/src/main/java/io/simplekafka/client/IdempotentProducer.java} 实现 \method{public ProduceReceipt send(TopicPartition tp,List<RecordData> records,Acks acks,long timeoutMillis)} 和 \method{public ProduceReceipt retry(TopicPartition tp)}；\method{IdempotentProducer(MetadataRouter router,long producerId,int producerEpoch)} 与 \method{close()} 已提供。每分区维护 nextSequence 与最多一个 pending batch；失败保留其冻结身份/数据，显式 retry 仅调用一次 RPC。
\end{enumerate}
\paragraph{错误、状态与副作用。}负 producer id/epoch/sequence、null topic partition/acks、null/empty batch、batch内 null、带 stamp 的输入、序列加法溢出、每条记录超上限、负 timeout 或超帧均为 \method{INVALID_REQUEST}，必须在追加前拒绝且不消耗sequence。新producer或更高producer epoch必须从sequence0开始；只有一条stamped记录持久化后新epoch才生效。当前epoch只接受精确nextSequence；旧epoch为 \method{FENCED_PRODUCER_EPOCH}，跳号/未记录旧序号为 \method{OUT_OF_ORDER_SEQUENCE}。同一已完成 producerId/producerEpoch/firstSequence 的 batchSize 不同为 \method{OUT_OF_ORDER_SEQUENCE}；batchSize 相同但 batchHash 不同为 \method{DUPLICATE_SEQUENCE_CONFLICT}。相同identity、边界和hash的完整重复返回原 \pathcode{[firstOffset,nextOffset)}，不追加、不改physical LEO；复制backend仍需检查当前leader epoch/minISR，并完整等待ALL的HW条件。若只有同身份尾batch前缀已持久化，精确原请求校验已有前缀后只写缺失suffix并返回整个原范围；其他producer/batch写入报 \method{INCOMPLETE_BATCH}，同身份不同payload报 \method{DUPLICATE_SEQUENCE_CONFLICT}。保留起点>0返回 \method{PRODUCER_STATE_EXPIRED}。
\paragraph{教学协议的单帧上限。}API 11 仍使用第11步的课程 TCP frame：length字段最小10、最大8,388,608 bytes（不含4-byte length prefix），最多8,388,598-byte payload，总帧最多8,388,612 bytes。producer须在发送前拒绝超过该帧限额的batch请求；记录自身仍受上项1,048,580-byte stamped record上限约束。
\paragraph{派生状态和调用边界。}producer-state cache由日志stamped记录推导。读取新尾部时增量扫描；prefixVersion改变或LEO回退时丢弃cache并从offset0重建；recover逐条执行 \pathcode{read(offset,1,1_048_580)}，检查每producer序列、batch index与batch hash。部分batch只允许留在日志尾部；recover可保留它等待原请求续写，完整但不合法的历史拒绝且不改原日志。存储mutation后遇到不确定I/O错误时不要把内存状态当成持久成功，failed log须重开恢复。leader \pathcode{epoch} 与 producer \pathcode{epoch} 不同：leadership切换后同一个producer身份仍可用原producer epoch重试。
\paragraph{client 边界。}\pathcode{IdempotentProducer(router,producerId,producerEpoch)} 构造时身份非负；负id/epoch为 \method{INVALID_REQUEST}。在可用producer上，\pathcode{send} 还拒绝null tp、records、acks，空批、null记录、负timeout、stamp输入、超记录/帧界或sequence溢出。metadata失败且数据RPC尚未发出时不构成未知追加。一次请求可能到达broker后失败，因此失败留下per-partition pending；同partition另一send为 \method{INVALID_REQUEST}，其他partition仍可发送。仅成功receipt且长度等于batch条数时推进nextSequence并清pending。无pending的 \pathcode{retry(tp)} 为 \method{INVALID_REQUEST}；retry先显式刷新metadata，再用原producer id/epoch、firstSequence、记录快照、acks、timeout向新route只发一次。刷新、RPC失败或非法receipt均保留pending。close清除client本地pending、不关闭借用router；已关闭client的send/retry抛 \pathcode{IllegalStateException}。
\lessoncontract{只对同一 producer identity、同一 batch边界和同一内容返回原持久范围；所有失败路径不得伪造receipt。复制错误不会回滚已追加的stamped记录。}
\lessonhints{先区分leader epoch、producer epoch、batch sequence、Kafka offset 和业务 eventId，各自不能互换。}{把持久日志看作真相、producer-state cache 看作带 prefixVersion 的派生结果；说明普通 append、retention 与 truncate 分别如何影响它。}{由存储开始按六项顺序实现并观察磁盘：先检查一条完整重复不改字节，再检查部分batch只补suffix，最后在API11丢回复后切换leader并由显式retry验证一份持久批次。}
\expected{假设 partition 无保留、leader broker1/epoch0、producerId=7、producerEpoch=0、firstSequence=0；batch是三条ASCII记录 \pathcode{("retry-0","value-zero",1001)}、\pathcode{("retry-1","value-one",1002)}、\pathcode{("retry-2","value-two",1003)}。key/value长度分别(7,10)、(7,9)、(7,9) bytes；普通记录分别49、48、48 bytes，stamped磁盘记录分别113、112、112 bytes。每条batchIndex为0/1/2，batchSize=3且共享同一hash。本例调用 \pathcode{send(tp,batch,Acks.LEADER,5_000)}；若API11回复丢失，客户端保留pending、leader物理LEO=3。复制并报告所有副本后HW=3；broker1停机、broker2成为leader/epoch1后，显式 \pathcode{retry(tp)} 保留原LEADER ACK与producerEpoch0/sequence0并返回原[0,3)，各日志仍三条。重新打开broker2磁盘，原请求仍由日志恢复识别并返回[0,3)，字节与LEO不变。若该重复请求使用 \pathcode{Acks.ALL}，broker2的ISR={2}<minISR=2，会在追加前报 \method{NOT_ENOUGH_REPLICAS}；\method{allAckDuplicateWaitsForHwAndStillRequiresMinIsrAndCurrentEpoch} 覆盖此边界。反例：若一个三条batch只写offset0–1，普通新请求得到 \method{INCOMPLETE_BATCH}；原身份/尺寸/hash/payload补交时仅写offset2，返回完整[0,3)。}
\paragraph{测试与完成标准。}\method{Step30Test.producerRejectsInvalidBatchesBeforeChangingSequenceOrLogs} 与 \method{sequenceErrorsDoNotAppendOrInstallAnUnpersistedEpoch} 检查输入、epoch/sequence错误无副作用；\method{stampedRecordDiskBytesAreIndependentAndStructurallyChecked} 检查磁盘格式/CRC/长度；\method{incompleteBatchSurvivesRecoveryAndCanOnlyBeCompletedByItsOriginalRequest}、\method{sameLeoRewriteCannotReuseOldProducerReceipt}、\method{recoveryRejectsCompleteBatchWithIncorrectFingerprint} 与 \method{retainedPrefixCannotServeCachedOrRecoveredProducerReceipts} 检查恢复、部分批次、前缀失效；\method{singleBrokerRejectsIdempotentProduceWithoutAppending} 与 \method{idempotentApiAndStampedFetchWireFormatsRoundTrip} 检查 API；\method{allAckDuplicateWaitsForHwAndStillRequiresMinIsrAndCurrentEpoch}、\method{lostResponseRetryAndDiskReopenPreserveOneDurableBatch}、\method{boundedReplicaFetchAndDiskReopenPreserveStampedSegments}、\method{reconciliationComparesProducerStampsAtCommittedAndUncommittedOffsets} 检查 ACK、真实TCP、磁盘重开与修复。单步命令定位Step30；累计命令运行Steps1–30，但第8章累计完成还需Step31。
\pitfalls{同一 producer retry 去重的是分区内带 sequence 的记录batch，不是不同offset上的业务事件。它不自动分配producer id，不是Kafka PID/session协议，不是transactional producer，也不保证业务副作用 exactly-once。普通 \method{SimpleProducer} 仍不重试；只有新client的调用方显式 \method{retry}。}
\lessoncommand{30}
\lessonquestions{为什么已完成重复batch仍须检查当前epoch、minISR与HW？}{retention删除日志前缀后，为什么不能猜回producer sequence或旧receipt？}
\paragraph{自检答案。}重复请求不需要新追加，但ALL承诺的是当前leader与副本已满足条件；复制确认条件可能随选举/ISR变化。被删前缀可能包含producer首次使用的epoch、sequence和原offset范围，没有这些持久证据不能安全区分新写与重复。

\lessonsection{31}{消费副作用幂等}
\goals{Kafka record 被重复投递时，仍只将每个稳定业务事件应用到账本一次。}{解释offset去重与eventId去重的差异，并实现可恢复的本地持久账本。}
\background{先修\stepref{30}{幂等生产者}，并回看\stepref{16}{处理后提交}与\stepref{29}{副作用重放}。Kafka offset 表示一条日志记录的位置，不是业务事件身份：同一事件可被重发到另一个offset；callback次数也可能因commit失败而增加。稳定 eventId 来自业务数据，若相同事件在重放和重开后保持相同ID，账本可以将它作为应用层去重键。这里实现的是一个本地强制落盘日志与内存余额的协调，不与Kafka offset store共享事务。}
\providededit{provided 层提供 \pathcode{BusinessEvent(String eventId,String account,long delta)}（\pathcode{provided/src/main/java/io/simplekafka/model/BusinessEvent.java}）与 \pathcode{MessageCodec.encodeBusinessEvent(BusinessEvent):byte[]} / \pathcode{decodeBusinessEvent(byte[]):BusinessEvent}（\pathcode{provided/src/main/java/io/simplekafka/protocol/MessageCodec.java}）。}{\pathcode{exercises/src/main/java/io/simplekafka/client/DurableEffectStore.java}：现有骨架的 \method{public DurableEffectStore(Path directory)} 构造器打开 \pathcode{PartitionLog} 并先调用recover；实现 \method{public synchronized boolean apply(BusinessEvent event)} 与 \method{public synchronized void recover()}。已有API为 \method{public synchronized long balance(String account)} 与 \method{public synchronized void close()}。}
\paragraph{业务事件与持久格式。}\pathcode{BusinessEvent} 的eventId/account都必须非空、严格UTF-8且各不超过255 bytes；delta为任意signed long，正数加款、负数减款、0也是合法。提供的value格式为 \pathcode{version:uint8=1 | eventId:length-prefixed identifier | account:length-prefixed identifier | delta:int64}；未知版本、截断、非法UTF-8/标识或尾随字节在codec层为 \method{INVALID_REQUEST}。BusinessEvent 构造时拒绝非法标识并抛 \pathcode{IllegalArgumentException}。store每个账本事件恰写一条 \pathcode{RecordData(null,encodedEvent,0)}，即key=null、timestamp=0、producerStamp=null；普通记录字节有32-byte固定开销。
\lessoncontract{\pathcode{apply} 对新合法eventId先用checked addition计算账户余额，再强制追加一条记录，最后发布seen与balance内存状态，成功返回true。完全相同的eventId/account/delta重放返回false，不追加、不改余额；同ID但account或delta不同为 \method{INVALID_REQUEST}，余额溢出也为 \method{INVALID_REQUEST}，两者都不追加。}
\begin{itemize}
\item 在store正常且日志完整时，\pathcode{apply(null)} 返回 \method{INVALID_REQUEST}；无效eventId/account无法形成合法 \pathcode{BusinessEvent}（提供构造器抛 \pathcode{IllegalArgumentException}）。\pathcode{balance(account)} 对不存在账户返回0；null、空、非法Unicode或超过255 UTF-8 bytes 的account返回 \method{INVALID_REQUEST}。
\item 同一store实例的apply/balance/recover/close同步。构造时和显式 \pathcode{recover()} 从offset0按序、逐次读取一条完整记录（\pathcode{maxRecords=1,maxBytes=1_048_580}）重建事件身份与余额；要求logStart=0，每记录key null、timestamp0、没有stamp。完整非法payload、格式字段不符、同ID历史冲突、余额累计溢出或scan无进度为 \method{CORRUPT_RECORD}，不得改写完整损坏日志。
\item 底层磁盘末尾只有不完整记录时，PartitionLog/SegmentLog恢复规则截到最后完整记录后再重建；完整但CRC/结构/业务编码错误不可按截断处理。若已存在重复的完全相同event历史，恢复只计第一次余额但保留日志中既有记录。
\item append I/O/返回范围不确定时不发布内存新余额并把实例标failed；截断账本(logStart>0)报 \method{CORRUPT_RECORD} 并failed。failed或closed之后apply/balance/recover为 \method{STORAGE_ERROR}；必须close后新建实例由磁盘恢复。close可重复调用。
\end{itemize}
\lessonhints{用业务eventId去重，不要使用Kafka offset；同一事件重投后offset可以不同但eventId仍相同。}{每个新事件先检查ID冲突和checked余额，再force追加；只有追加成功才更新内存seen/余额。}{先验证同一事件第二次apply不增长日志，再close并重开store重放历史，比较余额、seen行为和磁盘字节。}
\expected{设一个单分区输入日志按offset 0、1、2依序携带A、A、B，三条输入record都用null key、timestamp0、value为编码后的事件。A=\pathcode{BusinessEvent("event-A","acct",+10)}，重复A的eventId/account/delta完全相同；B=\pathcode{BusinessEvent("event-B","acct",-3)}。ASCII eventId分别7 bytes、account4 bytes；每个编码value为 \(1+4+7+4+4+8=28\) bytes，ledger普通记录占 \(32+28=60\) bytes。apply依次返回true、false、true；账本仅有[A,B]两条（120 bytes），余额7。若第一次callback处理全部三条后commit失败，callback history=[0,1,2]但committed仍空；新consumer从0重放并成功提交3，history=[0,1,2,0,1,2]，账本仍[A,B]且余额7，committed nextOffset=3。Kafka offset没有变成eventId，两个系统也没有共同事务。}
\paragraph{测试与完成标准。}\method{Step31Test.appliesAndRecoversDurableBusinessEffects} 检查首次/重复apply、余额和重开恢复；\method{rejectsConflictsInvalidIdentifiersAndOverflowWithoutAppending} 检查冲突、UTF-8边界与溢出；\method{businessEventCodecRejectsMalformedPayloads} 检查value wire格式；\method{concurrentDuplicateApplicationsAppendOnceAndRecoverOnce} 检查同步去重；\method{recoversRepeatedHistoryAndTruncatesIncompleteFinalRecord} 与 \method{rejectsCorruptCompleteHistoryWithoutChangingItsLog} 区分不完整尾与完整损坏；\method{consumerCommitFailureReplaysCallbacksButNotDurableBusinessEffects} 检查真TCP commit失败、6次callback、2条账本记录、余额7及committed=3。单步命令定位Step31；累计命令运行Steps1–31，是第8章完成边界。
\pitfalls{账本强制落盘只协调一个本地store的记录与内存发布；它不和producer、Kafka log或consumer offset形成分布式事务。相同事件可位于不同Kafka offset；stale group token能拒绝commit，却不能回滚账本。课程例子是应用层本地去重，不宣称全局transaction或端到端 exactly-once。}
\lessoncommand{31}
\lessonquestions{为什么稳定eventId比Kafka offset或callback次数适合作为业务去重键？}{账本追加成功但consumer offset commit失败后，哪些状态尚未协调？}
\paragraph{自检答案。}offset区分日志位置，callback计数只记录尝试；二者都不能识别跨位置的同一业务事件。账本已持久生效而commit仍未保存时，重放会重复调用callback，但eventId让账本返回false并避免重复余额变更；消费进度与账本仍不是原子提交。
